\section{Técnicas y decisiones descartadas}\label{sec:tecnicas-descartadas}

El análisis factorial (\cref{sec:af}) generó, además de su resultado final,
tres decisiones metodológicas que se probaron y se rechazaron de forma
explícita, y dos variables que el modelo no toleró. Documentarlas evita que
parezcan omisiones: cada una se descartó por una razón concreta y verificable,
no por conveniencia.

\subsection{Análisis paralelo de Horn}

Es el criterio estándar para decidir cuántos factores retener, y aquí retiene
siete --- el doble que el MAP de Velicer finalmente adoptado
(\cref{fig:af-paralelo}). No se sigue, porque el resultado es un artefacto del
tamaño muestral y no una señal sustantiva: el umbral simulado sobre la matriz
reducida se sitúa en 0.06 y sobre la completa en 1.06. Con las 16,008
observaciones de entrenamiento del bloque de 14 predictores, los eigenvalores que produciría ruido puro se concentran
fortísimo alrededor de su valor esperado, la banda de simulación se angosta
casi a un punto, y cualquier eigenvalor apenas positivo la supera. El
argumento para descartarlo se sostiene mirando los umbrales simulados, sin
necesidad de ver una sola carga factorial. El MAP de Velicer se usa en su
lugar por ser un mínimo descriptivo que no lleva $n$ en la fórmula y no se
infla con el tamaño de la muestra.

\subsection{Extracción por máxima verosimilitud}

La extracción de \cref{sec:af} usa ejes principales iterados y no máxima
verosimilitud, que exige normalidad multivariada. Las transformaciones
Box-Cox de \cref{sec:transformaciones} acercaron las marginales a la simetría
($|\text{sesgo}|<0.5$) pero no las volvieron normales, de modo que ese
supuesto no se sostiene sobre estos datos. La solución de ejes principales se
validó contra \texttt{psych::fa(fm="pa", rotate="varimax")} de R sobre el
mismo bloque: coincide hasta orden y signo de los factores, con comunalidades
iguales a la milésima.

\subsection{Rotación oblicua}

Antes de adoptar varimax se probó una rotación oblicua (oblimin), que permite
correlación entre factores. Encontró correlaciones de hasta 0.32 entre ellos
--- físicamente razonable, ya que el viento dispersa contaminantes y no es
independiente de la carga de combustión. Se rechaza para el uso instrumental
de los puntajes: el objetivo de \cref{sec:clasificacion_tecs} es disponer de
predictores no colineales, y una solución oblicua reintroduce exactamente el
problema de multicolinealidad que el análisis factorial vino a resolver. La
correlación entre factores se reporta como hallazgo, pero no se lleva al
modelo de clasificación.

\subsection{Dos variables que el modelo factorial no toleró}

\texttt{NOX} y \texttt{RH\_min\_diurna} produjeron comunalidades mayores que 1
--- los llamados \textbf{casos Heywood}, imposibles porque la comunalidad es
una proporción de varianza. \texttt{NOX} es, por definición química, la suma
de \texttt{NO} y \texttt{NO2} ($r=0.985$ contra esa suma): no es una variable
redundante en un sentido estadístico ordinario, es una identidad contable que
el PCA de \cref{sec:pca} absorbió sin protestar y que el análisis factorial no
puede repartir en comunalidad y unicidad. \texttt{RH\_min\_diurna}
correlaciona 0.892 con \texttt{RH\_media}; a diferencia de \texttt{NOX} no es
una identidad, así que \cref{sec:af} la dejó en observación hasta que el
segundo ajuste, ya con el diagnóstico completo, forzó la misma salida. Ambas
se eliminaron del bloque de predictores, que queda en 13 variables con un VIF
máximo de 2.54. Que el PCA no haya señalado ninguno de los dos problemas es en
sí mismo un resultado: el modelo factorial expone redundancias que la
descomposición de varianza total no tiene forma de reportar.
